\documentclass[11pt]{article}
\usepackage[utf8]{inputenc}
\usepackage[T1]{fontenc}
\usepackage{amsmath,amssymb}
\usepackage{array}
\usepackage{booktabs}
\usepackage{float}
\usepackage[margin=1in]{geometry}

\setlength{\emergencystretch}{3em}
\sloppy

\title{Review of docs/methodology.tex}
\author{documentation review pass}
\date{November 2026}

\begin{document}

\maketitle

\section{Scope and Method}

This review covers \texttt{docs/methodology.tex} (422 lines, compiled to 8 pages). A citation written
\texttt{tex:115} refers to a source line of that document. A citation naming
\texttt{settle\_expenses.py} with a line number refers to the implementation as it stands today (394
lines).

Evidence used:

\begin{itemize}
  \item Direct reading of \texttt{settle\_expenses.py} for every behavioural claim the document makes.
  \item \texttt{methodology.aux} for section numbering (sections 1--10, 8 pages).
  \item \texttt{methodology.log} for compile diagnostics (one warning, three overfull boxes).
\end{itemize}

Two runtime probes were not executed, so claims about executed behaviour below are reading-based, not
observed: the \texttt{\_settle()} selection question (finding C1) and the emptiness of the remainder
branch (finding M7). Both are argued from the source, and each is flagged as such.

\textbf{Note on process:} \texttt{methodology.tex} has not been modified. Everything below is a proposal
list; the file remains as authored.

\section{Findings at a Glance}

\begin{table}[H]
\centering
\begin{tabular}{>{\raggedright\arraybackslash}p{0.055\textwidth} >{\raggedright\arraybackslash}p{0.14\textwidth}
>{\raggedright\arraybackslash}p{0.515\textwidth} >{\raggedright\arraybackslash}p{0.16\textwidth}}
\toprule
ID & Location & Issue & Severity \\
\midrule
M1 & tex:115 & Theorem statement is false as written (feasible settlements are unbounded) & High \\
M2 & tex:126 & Lower-bound arithmetic error plus a missing step that the proof actually needs & High \\
M3 & tex:129, 222 & "Exactly $k-1$ payments" is false; counterexample with two payments of 5 & High \\
M4 & tex:32, 143 & $\min(n-c,\, n-1)$ is redundant; the second branch never occurs & Medium \\
M5 & tex:242 & Optimality proof cites a lemma that does not prove DP optimality & High \\
M6 & tex:209--220 & Reconstruction lemma rests on an unstated invariant & Medium \\
M7 & tex:224, 233 & "Otherwise $|P| \le n-c-1$" contradicts the same line's lower bound & High \\
M8 & tex:187 & Transition sum is off by $1/2$ (missing $w=0$ term) & Low \\
M9 & tex:336 & "Arbitrarily large" greedy gap: unsupported, and the described construction is impossible & Medium \\
M10 & tex:330 & "Exactly one heap is non-empty" at termination; both are empty & Medium \\
M11 & tex:334, 344 & Iteration bound $2n-1$ and "clears both removes one person" & Low \\
M12 & tex:370--371 & "Single participant has no effect" is false unless the payer participates & Medium \\
M13 & tex:376 & Share bound and "always positive" claim both wrong at $A=1, |U|=2$ & Medium \\
C1 & tex:139, 222, 398 & \texttt{\_settle()} pairs the smallest debtor, not the largest & High \\
C2 & tex:306--319 & Heap keys and pseudocode comment contradict the code and themselves & High \\
C3 & tex:49 & "Multiset of participants" contradicts duplicate rejection & Low \\
C4 & tex:393--401 & Source table header "Lines" and one wrong section mapping & Medium \\
C5 & tex:406--420 & Dispatch table omits \texttt{heuristic} and the unknown-method error & Low \\
C6 & tex:289--293 & Sorting cost contradicts its own "dominant term" line & Medium \\
L1 & whole file & No \texttt{\textbackslash label}; every hard-coded cross-reference is wrong & High \\
L2 & tex:281 & LaTeX commands inside \texttt{verbatim} print literally & Medium \\
L3 & tex:163, 213, 392--404 & Three overfull boxes; the source table overflows the text block & Medium \\
L4 & preamble & Seven loaded packages are never used & Low \\
\bottomrule
\end{tabular}
\end{table}

\section{Mathematical Findings}

\subsection{M1 --- The $n-1$ bound is stated backwards (tex:115)}

\begin{quote}
  Any feasible settlement uses at most $n-1$ payments.
\end{quote}

This is false: feasibility is closed under adding a redundant cycle of payments, so feasible settlements
are unbounded in size (take any settlement and add $x \to y$, $y \to z$, $z \to x$ of equal amounts).
The proof given is a construction, which establishes the existential form:

\begin{quote}
  There exists a feasible settlement with at most $n-1$ payments; hence the minimum size is at most $n-1$.
\end{quote}

The existential form is exactly what the later text needs (tex:233 uses it as an upper bound on the
constructed solution), so this is a wording fix, not a restructuring.

\subsection{M2 --- The $n-c$ lower bound proof does not prove the bound (tex:126)}

Two defects in one paragraph.

\textbf{Arithmetic.} The line reads
$\sum (k_j - 1) = n - \sum k_j$. With $\sum k_j$ denoting the total number of people inside the $c$
selected subsets, the left side equals $\sum k_j - c$, not $n - \sum k_j$. The two expressions coincide
only in the degenerate case where the document silently means something else.

\textbf{Direction.} Inside a lower-bound argument the remainder term is bounded from above ("the remaining
$n - \sum k_j$ people need at most ... additional payments"), which cannot yield a lower bound. Adding
the two halves as written gives $\ge \sum k_j - c$ plus an upper bound, i.e. nothing.

\textbf{What is missing.} The proof never says anything about the settlement being analysed. The standard
formal argument is:

\begin{enumerate}
  \item Take any feasible settlement $P$ with $t$ payments and build the multigraph with vertex set $S$
and one edge per payment.
  \item Each connected component has zero total balance (payments only transfer amounts inside a component),
so each component is a zero-sum subset of $S$.
  \item A component on $k$ vertices needs at least $k-1$ edges. Summing over components,
$t \ge \sum_{\text{comp}} (k-1)$, which equals $n$ minus the number of components.
  \item The components are disjoint zero-sum subsets, so their number is at most $c$. Hence $t \ge n - c$.
\end{enumerate}

Step 2 is the crux and is absent from the document. It is worth adding because it is also what makes the
result feel obvious once seen.

\subsection{M3 --- "Exactly $k-1$ payments" for greedy on a zero-sum group (tex:129, tex:222)}

tex:222 reads: "Since the subset is zero-sum, this terminates with all internal balances at zero in
exactly $k-1$ payments." A payment clears both endpoints whenever the two magnitudes are equal. Take
a zero-sum group of four people with balances $\{+5, +5, -5, -5\}$: two payments of 5 clear all four,
which is $k-2$, not $k-1$.

Correct statement: greedy uses \textbf{at most} $k-1$ payments; equality holds exactly when the payment
graph on the group is a tree. For the achievability half of Theorem 3.2 the honest construction is the
star: pick one person in the zero-sum subset as a hub, order the others, and pay each into or out of the
hub, giving exactly $k-1$. That proves existence without relying on any greedy behaviour.

\subsection{M4 --- $\min(n-c,\; n-1)$ has a branch that never fires (tex:32, tex:143, tex:129)}

Whenever $n \ge 1$ non-zero balances exist, $n \ge 2$ (balances sum to zero and none is zero) and the whole
set $S$ is itself a zero-sum subset, so $c \ge 1$ and $n - c \le n - 1$. The $\min$ in tex:32 and tex:143
is therefore always the first argument, and the parenthetical in tex:129 ("or $n-1$ if $n-c > n-1$,
i.e., $c < 1$") describes an unreachable case. Recommend stating the optimum as $n-c$ and, if a boundary
case is wanted, stating it explicitly as $n = 0$ (no payments).

\subsection{M5 --- The optimality proof is circular (tex:242)}

The implementation discussed here is \texttt{exact1}: bottom-up bitmask
dynamic programming with cached subset results. Its \texttt{dp} table
stores each subset's optimal group count for reuse by larger subsets.

\begin{quote}
  By Lemma 1, the DP finds the true maximum $c$.
\end{quote}

The only lemma in the document (tex:213) is "Correctness of reconstruction" and says nothing about
$\texttt{dp}$ computing a maximum. The load-bearing step is unproven. Add before it:

\begin{quote}
  \textbf{Lemma (DP value).} For every mask $m$, $\texttt{dp}[m]$ equals the maximum number of pairwise
disjoint zero-sum subsets contained in $m$.
  \textbf{Proof.} Induct on $|m|$. In an optimal packing for $m$, either the person of
$\texttt{lowbit}$ is unused --- value $\le \texttt{dp}[m \oplus \texttt{lowbit}]$ --- or they belong
to exactly one selected zero-sum subset $T \subseteq m$ containing that bit, giving
$1 + \texttt{dp}[m \setminus T] \le 1 + \texttt{dp}[m \oplus \texttt{candidate}]$ over candidates the
loop enumerates (tex:177 iterates every submask of $m \oplus \texttt{lowbit}$, so every such $T$ is
visited). Both branches match the transition, and both transitions are feasible.
\end{quote}

With that lemma, tex:242's argument becomes valid.

\subsection{M6 --- Reconstruction lemma has a gap that could hide a real bug (tex:209--220)}

The claim "The final $m$ with $\texttt{parent} = 0$ has no subsets remaining" does not follow from what
is proved. In the DP, $\texttt{parent}[m] = 0$ is simply the marker for Option A, and Option A can be
taken while $\texttt{dp}[m] > 0$ for a general mask. So the extraction loop
(\texttt{settle\_expenses.py:275--281}) could in principle stop early and return a partial packing
without any error surfacing.

The claim is nonetheless true along the reconstruction path, for a reason the document never states:
$m$ stays zero-sum throughout the loop. Initially $m = 2^n - 1$ is zero-sum; each extraction removes a
zero-sum subset. For a zero-sum $m$ with $\texttt{dp}[m] \ge 1$, take an optimal packing of $m$ that
leaves the $\texttt{lowbit}$ person uncovered: the leftover people are non-empty and sum to zero, so
they form one more disjoint zero-sum subset, contradicting maximality. Hence for such $m$ Option B strictly
beats Option A and $\texttt{parent}[m] \ne 0$. State that invariant (and the fact that it is what makes
the loop safe) in the proof.

\subsection{M7 --- The "remainder" branch contradicts the lower bound (tex:224, tex:233)}

tex:233 says, in one sentence: "Otherwise $|P| \le n - c - 1$. In all cases ... $|P| \ge n - c$". Those
two clauses cannot both hold. The resolution is that "otherwise" never happens, and the document should
say so:

\begin{quote}
  For any \textbf{maximal} packing of disjoint zero-sum subsets, the remainder is empty. The complement
of the packing sums to zero ($S$ sums to zero and each subset sums to zero); if it were non-empty it would
have at least two members (a single non-zero balance cannot sum to zero) and would be disjoint from every
selected subset, so it could be added --- contradicting maximality.
\end{quote}

That applies to both paths: the exact DP is maximum, and the heuristic is maximal because it scans every
enumerated zero-sum subset and takes any that is disjoint (\texttt{settle\_expenses.py:242--247}). So
\texttt{Step 5} (tex:224--226), the $|P_{\text{remainder}}|$ term in the tex:233 formula, and the
\texttt{if remaining:} branch (\texttt{settle\_expenses.py:153}) are all dead. Recommend either deleting
Step 5 and asserting $|P| = n - c$, or keeping the code branch and labelling it defensive.

\subsection{M8 --- Off-by-a-half in the transition count (tex:187)}

The chain
$\sum_{w=1}^{n} \binom{n}{w} 2^{w-1} = \frac12 \sum_{w=0}^{n} \binom{n}{w} 2^w = \frac{3^n}{2}$
silently includes the $w=0$ term, which is 1. The sum equals $(3^n - 1)/2$. The $O(3^n)$ conclusion
is unaffected; the middle equality is simply wrong as written.

\subsection{M9 --- "Arbitrarily large" greedy gap (tex:336)}

The sentence is unsupported and, read as a ratio, false. Greedy uses at most $n-1$ payments while the
optimum is $n-c$, and each zero-sum subset has at least two members so $c \le n/2$, giving
$\text{optimum} \ge n/2$. The ratio is therefore below 2 and the additive gap below $n/2$
(explicitly: gap $\le c - 1$). "Arbitrarily large" can only mean the additive gap as $n \to \infty$,
and the document gives no construction for that.

The parenthetical construction ("multiple disjoint zero-sum components connected through a single
\texttt{bridge} person") cannot be realised as described: a person with a non-zero balance cannot lie in
two disjoint zero-sum subsets of $S$, so there is no such bridge.

The repo does contain a concrete divergence, but it is heuristic-versus-exact, not greedy-versus-exact:
\texttt{tests/test\_heuristic\_vs\_exact.py} produces 6 payments with \texttt{exact} and 7 with
\texttt{heuristic} on the same input. Recommend replacing the sentence with that citable example, plus
the gap bound above.

\subsection{M10 --- Heap termination state (tex:330)}

"Exactly one heap is non-empty" is wrong: at termination both heaps are empty. Heap entries are only ever
pushed when a residual magnitude is strictly positive (\texttt{settle\_expenses.py:205--211} and
\texttt{216--220}), so a non-empty debtor heap implies a strictly negative total, contradicting
$\sum b_i = 0$. "At most one" (the earlier sentence) is fine; "exactly one" is not.

\subsection{M11 --- Iteration counting (tex:334, tex:344)}

tex:334: "Each iteration that fully clears both a debtor and a creditor removes one person from
consideration" --- clearing both removes two. tex:344: "Total iterations $\le 2n-1$" is loose. Since
$\texttt{amount} = \min(...)$ always exhausts at least one of the two people, at most $n-1$ iterations
run. The $O(n \log n)$ conclusion is unchanged either way.

\subsection{M12 --- "Single participant" is not a no-op (tex:370--371)}

\begin{quote}
  An expense with $|U| = 1$ credits the payer $A$ and debits the participant $A$ by $A$. ... No effect.
\end{quote}

This only holds when the payer \emph{is} the participant. If the lone participant is someone else, the
expense transfers an amount: the payer gets $+A$ and the participant $-A$, and \texttt{net\_balances} produces
exactly that. The next subsection (tex:373) already states the non-participant case, which is how the
error is visible: the two subsections contradict each other for $|U|=1$ and payer not in $U$.

\subsection{M13 --- Payer-as-participant bound and sign claim (tex:376)}

"Since $s_p \le A/m < A$ when $|U| \ge 2$, the payer always has a positive net from their own expense."
Two problems:

\begin{itemize}
  \item Shares are the rounded-up or rounded-down value, so the correct bound is
$s_p \le \lceil A/m \rceil$. That can equal $A$ even for $m \ge 2$.
  \item The conclusion fails when it does equal $A$: with $A = 1$, $m = 2$, and the payer among the
first $A \bmod m = 1$ participant, $s_p = 1 = A$ and the payer's net contribution from that expense is
$0$. Also $0$ for $A = 0$.
\end{itemize}

Correct claim: the payer's contribution is $A - s_p \ge 0$, strictly positive except when $A \le 1$
(with $m \ge 2$). Note also that the "Payer Not a Participant" subsection (tex:374) is correct as written
("the payer is a creditor for that expense" is a per-expense statement, which is the right framing).

\section{Document-versus-Code Findings}

\subsection{C1 --- \texttt{\_settle()} settles smallest-debtor-to-largest-creditor (tex:139, 222, 398)}

The document calls the intra-group routine "greedy largest-debtor-to-largest-creditor" three times
(tex:139, tex:222, tex:398). The implementation (\texttt{settle\_expenses.py:175}) is:

\begin{verbatim}
debtor = max(debtors, key=debtors.get)   # debtors values are all < 0
creditor = max(creditors, key=creditors.get)
\end{verbatim}

\texttt{debtors} is built from balances $< 0$ (\texttt{settle\_expenses.py:170}), so
$\texttt{max}$ over it selects the \textbf{least negative} balance, i.e. the person who owes the
\emph{least}. The creditor side is correct (largest positive). So the executed intra-group rule is
smallest-debtor-to-largest-creditor, and the function's own docstring (\texttt{settle\_expenses.py:161})
repeats the same incorrect description.

Consequences worth stating in the review:

\begin{itemize}
  \item Section 4.5 (tex:220--222) and the Source Reference role text (tex:398) describe behaviour the
code does not have.
  \item The abstract's bullet on \texttt{greedy} (tex:35) and Section 6 describe the two-heap routine,
which \emph{is} largest-to-largest (\texttt{settle\_expenses.py:199--200}). Two different greedy rules
therefore live in the module, and the document describes one of them twice.
  \item Tests that assert payment counts and zero-sum validity would not catch this; only the tie-breaking
and pairing pattern would differ.
\end{itemize}

This is the item most likely to be a genuine code bug rather than a document bug, and it needs a
decision from the maintainer: either \texttt{max(debtors, ...)} becomes
\texttt{min(debtors, key=debtors.get)} (or \texttt{max(..., key=abs)}) to match the documented rule, or
the doc and docstring change to describe the current rule. Reading-based: not confirmed by execution.

\subsection{C2 --- Heap keying and a self-contradictory comment (tex:306--307, tex:314--319)}

The document says both heaps are keyed by $(-b_i, i)$. Only the creditor heap is:

\begin{verbatim}
debtors   = [(b, p)  for p, b in balances.items() if b < 0]   # key is b itself
creditors = [(-b, p) for p, b in balances.items() if b > 0]   # key is -b
\end{verbatim}
(\texttt{settle\_expenses.py:199--200})

Debtors are already negative, so their raw balance is the min-heap key for largest magnitude. The document's
pseudocode then carries the comment \texttt{k\_d = -b\_d < 0}, which contradicts itself: if $b_d < 0$ then
$-b_d > 0$. With that comment taken literally, the following line
(\texttt{amount = min(-k\_d, -k\_c)}) would compute a negative amount. Taking the code's convention
($k_d = b_d < 0$) makes the formula right. Recommend fixing the comment and splitting the two bullets
so the asymmetry is explicit. Also worth one sentence on tie-breaking: the heap path breaks ties
lexicographically on the person key, while \texttt{\_settle()} breaks them by dict insertion order.

\subsection{C3 --- "Multiset of participants" (tex:49)}

\texttt{net\_balances} would debit a repeated name twice, but \texttt{validate\_expense} rejects
duplicates inside one expense (\texttt{settle\_expenses.py:330}). The document's input contract (Section 1)
never mentions validation or the no-duplicates rule, so "multiset" is both wrong and the only place a
reader would infer it. Suggest "list of distinct participants, in CSV order" --- CSV order is load-bearing
for the rounding rule (tex:68--78), which the input list should flag.

\subsection{C4 --- Source Reference table (tex:390--404)}

Three issues in one table:

\begin{itemize}
  \item The header is "Lines" but the entries are \texttt{\S2}, \texttt{\S3}, ... which render as section
signs pointing at this document's own sections, not source lines of \texttt{settle\_expenses.py}. Rename
the column or replace the entries with real line ranges.
  \item \texttt{settle\_expenses()} $\to$ \texttt{\S3} is off: Section 3 is the graph formulation. Its
dispatch behaviour is Section 10 (tex:406), and its strategy is Section 4.1 (tex:134). The other five
mappings check out (\texttt{net\_balances} 2, \texttt{\_settle} 4, \texttt{\_greedy\_settle\_balances} 6,
exact DP 4, heuristic 5).
  \item \texttt{validate\_expense()} and \texttt{parse\_csv()} are absent, while the section's intro claims
the table covers the module's functions. Two rows would close it.
\end{itemize}

\subsection{C5 --- Dispatch table is incomplete (tex:406--420)}

The four cases listed cover \texttt{exact}, \texttt{greedy}, and both \texttt{auto} branches. Missing:
\texttt{method="heuristic"} (accepted, never warns, no bound applied, and the only path that is neither
proven minimal nor polynomial), and the \texttt{ValueError} for an unrecognised method
(\texttt{settle\_expenses.py:98--101}). Nothing in the section is stale; the signature at
\texttt{settle\_expenses.py:74} is (abridged, type hints omitted)

\begin{verbatim}
def settle_expenses(expenses, method="auto", threshold=18)
\end{verbatim}

and the fallback path emits a \texttt{UserWarning} (\texttt{settle\_expenses.py:110--114}).

\subsection{C6 --- Heuristic sorting cost contradicts the "dominant term" line (tex:289--293)}

"Sorting: $O(n^2 \cdot 2^n)$" followed by "Dominant term: $O(n \cdot 2^n)$" cannot both be right. The
implementation sorts by subset length (\texttt{settle\_expenses.py:239}), i.e. $O(z \log z)$ for $z$
zero-sum subsets with $z \le 2^n$, which is $O(n \cdot 2^n)$ --- consistent with the dominant term and with
the comparison table's $O(n \cdot 2^n)$ (tex:353). The space claim $O(n \cdot 2^n)$ (tex:355) is right
but should say why: the module keeps every zero-sum subset in memory as frozensets
(\texttt{settle\_expenses.py:232--237}).

\section{LaTeX and Presentation Findings}

\subsection{L1 --- No cross-references resolve (whole file)}

There is no \texttt{\textbackslash label} anywhere in the document, and every reference is hard-coded.
Because the preamble declares section-based theorem counters
(\texttt{\textbackslash newtheorem\{theorem\}\{Theorem\}[section]} with \texttt{lemma} sharing the same
counter), the environments actually print as:

\begin{itemize}
  \item Theorem 2.1 (Zero-Sum Invariant),
  \item Theorem 3.1 (the $n-1$ bound), Theorem 3.2 (the $n-c$ bound),
  \item Lemma 4.1 (Correctness of reconstruction), Theorem 4.2 (Optimality).
\end{itemize}

The prose cites "Theorem 2" (tex:156, tex:330), "Theorem 3" and "Theorem 4" (tex:233), "Theorem 4"
again and "Lemma 1" (tex:242). None match, and "Theorem 4" does not exist. A shared counter also means a
reader cannot tell a lemma from a theorem by its number. Fix: add \texttt{\textbackslash label}s and
\texttt{\textbackslash ref}s (plus \texttt{hyperref}) and let the numbering settle itself.

\subsection{L2 --- Markup inside \texttt{verbatim} prints literally (tex:281)}

The heuristic pseudocode block contains \texttt{if subset \textbackslash cap used == \textbackslash emptyset:}
inside \texttt{verbatim}. That prints as literal backslash markup in the PDF rather than as set
notation. Either spell it in plain ASCII (\texttt{if not (subset \& used):}, matching the code at
\texttt{settle\_expenses.py:244}) or move the line out of \texttt{verbatim}.

\subsection{L3 --- Overfull boxes (from \texttt{methodology.log})}

\begin{verbatim}
Overfull hbox (147.70341pt too wide) detected at line 163
Overfull hbox (7.29892pt too wide) in paragraph at lines 213-214
Overfull hbox (72.94775pt too wide) in paragraph at lines 392-404
LaTeX Warning: No \author given.
\end{verbatim}

The 147.7pt case is the \texttt{dp[m]} display at tex:163, whose single-line \texttt{\textbackslash text\{...\}}
prose is far wider than the text block; break it into two lines or state it in prose with an inline
formula. The 72.9pt case is the Source Reference table (tex:392--404): it visibly exceeds the 1in margins,
so it needs narrower column types (\texttt{p\{...\}}) or short roles. The 7.3pt one is cosmetic.
Adding an \texttt{\textbackslash author} removes the warning.

\subsection{L4 --- Unused packages (tex:5--18)}

\texttt{longtable}, \texttt{multirow} and \texttt{makecell} are loaded and never used (no matching
commands appear in the body), and \texttt{array}, \texttt{colortbl}, \texttt{xcolor} and \texttt{amsfonts}
likewise have no user: no coloured cells, no custom column specs, no \texttt{\textbackslash mathbb}. Only
\texttt{amsmath}/\texttt{amssymb} (math and \texttt{\textbackslash checkmark}), \texttt{amsthm},
\texttt{booktabs}, \texttt{float} (two \texttt{[H]} tables) and \texttt{geometry} are doing work.

\subsection{L5 --- Abstract overclaims (tex:38)}

"We prove correctness for all three" is not met: the exact path's proof is circular (M5), the greedy proof
has the heap-state error (M10), and the heuristic section offers no proof at all --- only the assertion that
it "may" miss the optimum (tex:295--296), with no example in the document. One worked example there, or
a pointer to the test file, would make the claim earnable.

\section{What Checks Out}

Worth recording so the fixes above can be targeted rather than sweeping:

\begin{itemize}
  \item The rounding rule (tex:68--82) matches \texttt{divmod} in \texttt{net\_balances}
(\texttt{settle\_expenses.py:57--66}) exactly, including "the remainder allocated one at a time to earlier participants in CSV
order", and the partition identity $\sum_u s_u = A$ is correct.
  \item The Zero-Sum Invariant (tex:84--96) and its proof are correct and are used correctly later.
  \item The bitmask encoding, the zero-sum precomputation cost $O(n \cdot 2^n)$, the submask iteration
(tex:177, matching \texttt{settle\_expenses.py:263--272}), the Option A/Option B dichotomy, and the
$O(3^n)$/$O(2^n)$ totals are all right (modulo M8's off-by-a-half).
  \item The threshold rationale (tex:263--265) matches the code's default of 18 (\texttt{settle\_expenses.py:74}).
  \item Payment-amount positivity and the no-self-payment invariant (tex:379--384) hold in both routines.
  \item The comparison table's time/space row (tex:352--355) is consistent with the code once C6's sorting
line is fixed.
\end{itemize}

\section{Suggested Order of Repair}

\begin{enumerate}
  \item \textbf{C1} --- decide whether \texttt{\_settle()} should pick the largest or smallest debtor, then
align the code, its docstring, and tex:139/222/398. Everything else in Section 4 depends on the answer.
  \item \textbf{M5 + M2 + M3} --- restore a real chain: DP-value lemma, component-graph lower bound, star
construction for achievability.
  \item \textbf{M7} --- state that the remainder is always empty; delete or label Step 5.
  \item \textbf{M1, M6, M10, M12, M13} --- wording and boundary corrections with no structural cost.
  \item \textbf{L1} --- switch to \texttt{label}/\texttt{ref}; until then every reference in the file reads
wrong.
  \item \textbf{C2, C4, C5, C6, M9, M11, M8, L2--L5} --- local corrections.
\end{enumerate}

\subsection{Claims not verified by execution}

\begin{itemize}
  \item C1 is a reading of \texttt{settle\_expenses.py:175}. A probe of \texttt{\_settle()} for four
people with balances \texttt{\{A:5, B:5, C:-5, D:-5\}} was declined at the shell, so the executed pairing
pattern is inferred, not observed.
  \item M7's "remainder is always empty" is a proof, not a survey of runs; if any path is found where a
remainder appears, the maximality argument above is what would need re-examining.
\end{itemize}

\end{document}
